% Generated draft: requires mathematical and source review.
\documentclass[11pt]{article}
\usepackage[margin=25mm]{geometry}
\usepackage{fontspec}
\setmainfont{Latin Modern Roman}
\usepackage{amsmath,amssymb,mathtools,unicode-math}
\setmathfont{Latin Modern Math}
\usepackage{enumitem,needspace,xurl,hyperref}
\setlength{\parindent}{0pt}
\setlength{\parskip}{5pt}
\setlength{\emergencystretch}{3em}
\newcommand{\N}{\mathbb{N}}
\newcommand{\Z}{\mathbb{Z}}
\newcommand{\Q}{\mathbb{Q}}
\newcommand{\R}{\mathbb{R}}
\newcommand{\C}{\mathbb{C}}
\newcommand{\F}{\mathbb{F}}
\newcommand{\A}{\mathbb{A}}
\newcommand{\PP}{\mathbb{P}}
\newcommand{\E}{\mathbb{E}}
\newcommand{\eps}{\varepsilon}
\newcommand{\dd}{\,\mathrm{d}}
\begin{document}
\section*{UnsolvedMath ID: 30006599}
\textbf{Source problem code:} catalogue entry
\subsection{Definitions and mathematical statement}
Let \(L\) be a finite lattice, with least element \(\hat 0\), greatest element \(\hat 1\), join \(\vee\), and meet \(\wedge\). The lattice is \emph{semidistributive} if, for all \(x,y,z\in L\),
\[
 x\vee y=x\vee z\ \Longrightarrow\ x\vee(y\wedge z)=x\vee y,
 \qquad
 x\wedge y=x\wedge z\ \Longrightarrow\ x\wedge(y\vee z)=x\wedge y.
\]
An element \(j\ne\hat 0\) is join-irreducible if it has a unique lower cover, denoted \(j_*\). For a finite semidistributive lattice, write \(C(x)\) for the canonical join representation of \(x\): the unique irredundant set of join-irreducibles whose join is \(x\) and which refines every join representation of \(x\). Here a set \(A\) refines \(B\) if every \(a\in A\) satisfies \(a\le b\) for some \(b\in B\). Define
\[
 \kappa(j)=\max\{y\in L:j\wedge y=j_*\},
 \qquad
 R_L(x)=\bigwedge_{j\in C(x)}\kappa(j).
\]
These maxima exist in this setting, and \(R_L\) is a rowmotion convention. Empty joins and meets are interpreted as \(\hat 0\) and \(\hat 1\), respectively. The supplied record does not define its rowmotion convention; it may use this bijection or its inverse. Consequently, agreement with the source's particular direction remains to be checked.

\textbf{Requested construction.} For every finite semidistributive lattice \(L\), find a homological interpretation of its rowmotion bijection, extending the Auslander--Reiten interpretation that the supplied record reports for finite distributive lattices. In a convention using \(R_L\), the induced permutation should agree with \(R_L\); the inverse convention requires the corresponding adjustment. The record does not specify the category, algebra, homological operation, or identification with lattice elements that such an interpretation must use, nor what precise compatibility constitutes extending the distributive case. Thus this is a request for a construction with unresolved specifications, rather than a fully specified existence theorem.

\textbf{Supplied contextual excerpt.} For a reverse linear extension \(\pi\), the elements of \(L\) are numbered from top to bottom. The excerpt includes
\[
 \Phi_\pi(\dim P_i)=-\dim I_i.
\]
It calls a matrix \(PU\) if it has the form \(wU\), with \(w\) a permutation matrix and \(U\) upper triangular, and describes \(w\) by the columns containing the first nonzero entries of the rows. Its Theorem 2, attributed only to \([\mathrm{KMT}]\), states that a lattice \(L\) is Auslander regular if and only if \(\Phi_\pi\) is \(PU\) for every reverse linear extension \(\pi\), equivalently for some such \(\pi\); in that case it asserts \(w_\pi=\rho_\pi\) when \(\Phi_\pi\in w_\pi U\). These are supplied contextual claims, not additional established hypotheses or a solution. The excerpt does not define the relevant algebra, \(\Phi_\pi\), \(P_i\), \(I_i\), or \(\rho_\pi\). Moreover, the stated first-nonzero-entry characterization and uniqueness of \(w\) require qualifications absent from the bare factorization with an arbitrary upper-triangular matrix, such as nonsingularity of \(U\).
\subsection{Short English statement}
For every finite semidistributive lattice, seek a homological construction realizing rowmotion and recovering the supplied Auslander--Reiten interpretation in the distributive case. The source leaves the intended homological setting, compatibility requirement, and direction convention unspecified. Its open status and verification claims are supplied curation, not independently verified facts.
\subsection{Sources}
\begin{itemize}
\item Local UnsolvedMath registry, \texttt{data/unsolvedmath/problems.json}, permanent ID \(30006599\), source code \texttt{catalogue entry}, title: \emph{Homological Interpretation of Semidistributive Rowmotion}. The supplied record attributes the problem to the Oberwolfach Reports collection and gives proposed year \(2026\); no individual proposer is supplied.
\item \emph{Representation Theory of Quivers and Finite-Dimensional Algebras} (2026), Oberwolfach Reports, DOI \(10.4171/owr/2026/6\), pp.\ 345--432. Supplied source and literature link: \url{https://doi.org/10.4171/owr/2026/6}. The supplied assessment identifies the problem as Question 1 and reports that the distributive Auslander--Reiten interpretation leaves the semidistributive, non-distributive case unresolved.
\item Defant et al., \emph{Rowmotion and echelonmotion}. Supplied literature link: \url{https://arxiv.org/abs/2507.18230}.
\item The contextual excerpt attributes its Theorem 2 to \([\mathrm{KMT}]\), without supplying a full bibliographic reference.
\item Literature assessment date: \texttt{2026-08-22}. The supplied assessment labels the problem open and describes its review as dated literature triage derived from supplied Oberwolfach curation, with time-sensitive status and cited claims requiring independent verification. This date and assessment are supplied dated curation, not independently verified here.
\item The registry labels the statement \texttt{corrected\_verified} and says it was reconstructed and checked against the official report. This is a supplied source-verification assertion; no source visit or independent verification is claimed here.
\end{itemize}
\end{document}
